\section{Kiểm thử và đánh giá hiệu năng}
\label{ch:testing}

\subsection{Kiểm thử chức năng và hợp đồng}

Lượt nghiệm thu sau cleanup ngày 10/11/2026 (\path{tasks/PHASE_7_9_FINAL_ACCEPTANCE_PLAN.md}) ghi nhận:

\begin{itemize}
  \item Toàn bộ backend solution: \textbf{285/285 test pass}.
  \item Contract validation: \textbf{13/13 pass}; QoS failure-policy: \textbf{8/8 pass}.
  \item Runtime acceptance cho baseline, đặt lịch, tranh chấp slot, broker gián đoạn và phục hồi, soak, queue/medical record và độ trễ event.
\end{itemize}

Bộ test hiện có 267 phương thức \code{[Fact]}/\code{[Theory]} trải trên 12 project test; số test case lớn hơn số phương thức vì mỗi \code{[Theory]} chạy với nhiều bộ dữ liệu.

\subsection{Phương pháp kiểm thử tải}

Kiểm thử tải dùng k6 với các kịch bản trong \path{tests/load/}: smoke, load, spike, soak, contention, booking, booking high-load và mixed workload. Nguyên tắc so sánh:

\begin{codeitems}
  \item Hai stack có cùng tổng giới hạn 12,5 CPU, chạy lần lượt để không tranh tài nguyên.
  \item Dùng cùng script cho cả hai hệ thống; mỗi lượt kiểm tra trạng thái nghiệp vụ (Confirmed, idempotency) và dọn dữ liệu test.
  \item \code{http\_req\_failed} của k6 đếm cả 409 tranh chấp hợp lệ, nên tỷ lệ lỗi server được tính riêng từ 5xx.
  \item Evidence chỉ lưu metrics tổng hợp; tài khoản test dùng file env bên ngoài repo.
\end{codeitems}

\subsection{Benchmark đặt lịch 20 VU}

Lượt đầu ngày 11/11 cho p95 booking API 594,7 ms ở microservices và 293,4 ms ở monolith. Sau khi tối ưu truy vấn reserve (mục~\ref{sec:optim}), ngày 17/11 đo lại mỗi bên ba lượt (Bảng~\ref{tab:bench20}, Hình~\ref{fig:bench20}).

\begin{table*}[t]
\centering
\caption{Benchmark đặt lịch 20 VU, 12,5 CPU mỗi hệ thống (17/11/2026)}
\label{tab:bench20}
\small
\begin{tabularx}{\linewidth}{@{}lRRRl@{}}
\toprule
\textbf{Hệ thống, lượt} & \textbf{Booking API p95} & \textbf{POST~$\rightarrow$~Confirmed p95} & \textbf{Thành công} & \textbf{p95 < 500 ms} \\
\midrule
Microservices, lượt 1 & 983,8 ms & 1.529 ms & 20/20 & Không đạt \\
Microservices, lượt 2 & 417,4 ms & 729,6 ms & 20/20 & Đạt \\
Microservices, lượt 3 & 449,0 ms & 661,1 ms & 20/20 & Đạt \\
\midrule
Monolith, lượt 1 & 796,7 ms & 908,2 ms & 20/20 & Không đạt \\
Monolith, lượt 2 & 335,6 ms & 408,0 ms & 20/20 & Đạt \\
Monolith, lượt 3 & 297,9 ms & 368,6 ms & 20/20 & Đạt \\
\bottomrule
\end{tabularx}
\end{table*}

\begin{figure}[!t]
\centering
\begin{tikzpicture}
\begin{axis}[
  ybar, bar width=9pt, width=0.95\linewidth, height=5.4cm,
  ylabel={Booking API p95 (ms)}, ymin=0, ymax=1100,
  symbolic x coords={r1,r2,r3}, xtick=data,
  xticklabels={Lượt 1,Lượt 2,Lượt 3},
  enlarge x limits=0.25,
  legend style={at={(0.5,-0.22)},anchor=north,font=\small,legend columns=2,/tikz/every even column/.append style={column sep=8pt}},
  nodes near coords, nodes near coords style={font=\scriptsize},
  every node near coord/.append style={/pgf/number format/fixed,/pgf/number format/precision=0},
  extra y ticks={500}, extra y tick style={grid=major,grid style={red!60,dashed,thick}},
  extra y tick labels={}
]
  \addplot[fill=micro!70,draw=micro] coordinates {(r1,983.8) (r2,417.4) (r3,449.0)};
  \addplot[fill=mono!60,draw=mono]  coordinates {(r1,796.7) (r2,335.6) (r3,297.9)};
  \legend{Microservices, Monolith}
\end{axis}
\end{tikzpicture}
\caption{p95 booking API qua ba lượt đo; đường đứt đỏ là ngưỡng 500 ms}
\label{fig:bench20}
\end{figure}

Mọi lượt đạt 101/101 check, cleanup đầy đủ và không có 5xx. Lượt đầu của cả hai bên chậm do khởi động nguội. Median p95 của microservices giảm khoảng 24,5\% so với mốc 594,7 ms nhưng vẫn cao hơn monolith khoảng 34\% (449,0 so với 335,6 ms).

\subsection{Booking high-load 30 VU, 2.400 slot}

Kịch bản: 30 VU, 10 giây ramp-up, 60 giây giữ tải, 10 giây ramp-down, tối đa 2.400 slot (10 bác sĩ $\times$ 5 ngày). Mỗi hệ thống chạy ba lượt acceptance.

\begin{table}[!t]
\centering
\caption{Booking high-load, median của ba lượt}
\label{tab:highload}
\small
\begin{tabularx}{\linewidth}{@{}L{3.1cm}RR@{}}
\toprule
\textbf{Chỉ số} & \textbf{Monolith (22/11)} & \textbf{Microservices (17/11)} \\
\midrule
p95 mọi booking API & 190,2 ms & 253,6 ms \\
p95 API tạo booking thành công & 214,6 ms & 387,8 ms \\
p95 POST~$\rightarrow$~Confirmed & 235 ms & 762 ms \\
Thông lượng booking Confirmed & 29,84/s & $\approx$ 29,72/s \\
Tỷ lệ 5xx & 0\% & 0\% \\
\bottomrule
\end{tabularx}
\end{table}

\begin{figure}[!t]
\centering
\begin{tikzpicture}
\begin{axis}[
  xbar, bar width=7pt, width=0.72\linewidth, height=5.2cm,
  xlabel={p95 (ms)}, xmin=0, xmax=900,
  symbolic y coords={e2e,create,all},
  ytick=data,
  yticklabels={POST $\rightarrow$ Confirmed,API tạo thành công,Mọi booking API},
  y tick label style={font=\small},
  enlarge y limits=0.3,
  legend style={at={(0.5,-0.32)},anchor=north,font=\small,legend columns=2,/tikz/every even column/.append style={column sep=8pt}},
  nodes near coords, nodes near coords style={font=\scriptsize}
]
  \addplot[fill=micro!70,draw=micro] coordinates {(762,e2e) (387.8,create) (253.6,all)};
  \addplot[fill=mono!60,draw=mono]  coordinates {(235,e2e) (214.6,create) (190.2,all)};
  \legend{Microservices, Monolith}
\end{axis}
\end{tikzpicture}
\caption{So sánh p95 trong kịch bản booking high-load}
\label{fig:highload}
\end{figure}

Cả hai đạt mọi threshold và dữ liệu test được dọn sạch. Thông lượng xác nhận gần bằng nhau, còn độ trễ của monolith thấp hơn. Khoảng cách lớn nhất nằm ở POST~$\rightarrow$~Confirmed, vì microservices xác nhận bất đồng bộ qua worker và một lời gọi HTTP sang Scheduling, còn monolith xác nhận trong cùng transaction.

Cần lưu ý: monolith dùng chung DB qua các lượt nên lịch sử appointment đã hủy tích lũy, và p95 tăng theo thứ tự chạy (122 $\rightarrow$ 190 $\rightarrow$ 254 ms). Một lượt monolith bị loại khỏi nhóm acceptance do teardown timeout, dù các threshold nghiệp vụ đều đạt.

\subsection{Mixed workload}

Mixed workload mô phỏng nhiều persona cùng lúc. Profile \textit{normal} gồm tìm kiếm và đặt lịch; \textit{clinic-peak} thêm lễ tân tiếp nhận walk-in, check-in và toàn bộ vòng đời hàng đợi.

\begin{table*}[t]
\centering
\caption{Mixed workload, mỗi profile một lượt cho mỗi hệ thống}
\label{tab:mixed}
\footnotesize
\begin{tabularx}{\linewidth}{@{}L{2.6cm}L{2cm}L{2.2cm}YL{2.1cm}L{1.9cm}@{}}
\toprule
\textbf{Profile} & \textbf{Hệ thống} & \textbf{Đọc / ghi} & \textbf{p95 đọc / ghi / tạo booking} & \textbf{Thành công / 5xx} & \textbf{Queue} \\
\midrule
\multirow{2}{2.6cm}{normal, 20 VU, 5 phút} & Monolith & 68,75 / 31,25\% & 82,2 / 111,3 / 50,3 ms & 100\% / 0\% & — \\
 & Microservices & 71,30 / 28,70\% & 26,8 / 91,3 / 107,2 ms & 100\% / 0\% & — \\
\midrule
\multirow{2}{2.6cm}{clinic-peak, 30 VU, 5 phút} & Monolith & 64,00 / 36,00\% & 276,7 / 362,4 / 287,8 ms & 94,81\% / 2,90\% & 0/865 \\
 & Microservices & 67,73 / 32,27\% & 49,7 / 187,8 / 343,8 ms & 99,976\% / 0\% & 836/836 \\
\bottomrule
\end{tabularx}
\end{table*}

Ở profile normal, microservices có p95 đọc và ghi thấp hơn, còn monolith có p95 tạo booking thấp hơn. Clinic-peak trên monolith \textbf{không đạt} do một lỗi chức năng: luồng walk-in lấy \code{WorkSchedule} bằng \code{AsNoTracking} rồi thêm \code{Appointment} qua navigation graph, khiến EF Core cố insert lại schedule đã tồn tại và ném \code{DbUpdateException}. Vì vậy chưa thể so hiệu năng hàng đợi giữa hai hệ thống.

\subsection{Tối ưu đã thực hiện}
\label{sec:optim}

\begin{codeitems}
  \item \textbf{Truy vấn kiểm tra capacity khi reserve} (\path{SchedulingManager.cs}): trước đây đọc toàn bộ \code{TimeSlot} của các reservation còn hiệu lực về ứng dụng; nay dùng một truy vấn aggregate theo \code{WorkScheduleId}, chỉ trả về tổng reservation và số lần slot yêu cầu xuất hiện. Khóa bản ghi schedule và filtered unique index vẫn giữ nguyên.
  \item \textbf{Đánh thức worker theo tín hiệu} (\path{AppointmentWorkSignal.cs}): outbox publisher và confirmation worker được đánh thức ngay khi có việc mới thay vì chờ chu kỳ polling, giảm thời gian POST~$\rightarrow$~Confirmed.
  \item \textbf{Connection pool:} cấu hình \code{Min Pool Size}/\code{Max Pool Size} cho Scheduling, Appointment, Queue và Identity.
\end{codeitems}

\subsection{Nhận định so sánh}

Với các lượt đo hiện có:
\begin{itemize}
  \item Monolith có độ trễ đặt lịch thấp hơn, nhờ ít hop mạng và xác nhận trong một transaction cục bộ.
  \item Thông lượng xác nhận booking của hai hệ thống tương đương (khoảng 29,7--29,8 booking/giây).
  \item Microservices xử lý được tải hỗn hợp có hàng đợi, trong khi monolith gặp lỗi chức năng ở luồng walk-in.
\end{itemize}

Các kết quả này \textbf{chưa đủ để kết luận kiến trúc nào tốt hơn tổng thể}, vì chưa khôi phục DB snapshot giữa các lượt, chưa thu CPU/RAM của hai stack theo cùng cách, và mixed workload mới có một lượt mỗi bên.
